% ============================================================================
% SEC03.TEX - Section 5.3: Mekanisme Menembak
% ============================================================================

\section{Mekanisme Menembak}

Pesawat tempur tidak lengkap tanpa senjata. Kita akan membuat pesawat menembakkan prefab \texttt{Bullet} yang sudah kita buat.

\subsection{Update PlayerController.cs}

Tambahkan variabel dan logika menembak:

\begin{lstlisting}[language={[Sharp]C}]
public GameObject bulletPrefab;
public Transform firePoint;

void Update()
{
    // ... input movement ...

    // Deteksi tombol Spasi
    if (Input.GetKeyDown(KeyCode.Space))
    {
        Shoot();
    }
}

void Shoot()
{
    Instantiate(bulletPrefab, firePoint.position, firePoint.rotation);
}
\end{lstlisting}

\subsection{Mengatur Fire Point}
Objek peluru harus muncul dari moncong pesawat, bukan dari tengah.
\begin{enumerate}
    \item Buat Empty GameObject baru sebagai child dari \texttt{Player}, beri nama \texttt{FirePoint}.
    \item Geser posisinya sedikit ke depan pesawat (sumbu Y positif).
    \item Di Inspector \texttt{Player}, drag prefab \texttt{Bullet} ke slot \texttt{Bullet Prefab}.
    \item Drag object \texttt{FirePoint} ke slot \texttt{Fire Point}.
\end{enumerate}

\subsection{Script BulletController}
Peluru perlu bergerak maju sendiri setelah muncul. Buat script \texttt{BulletController.cs} dan pasang pada prefab \texttt{Bullet}:

\begin{lstlisting}[language={[Sharp]C}]
public class BulletController : MonoBehaviour
{
    public float speed = 10f;

    void Update()
    {
        transform.Translate(Vector2.up * speed * Time.deltaTime);
        
        // Hapus peluru jika sudah jauh (setelah 2 detik)
        Destroy(gameObject, 2f); 
    }
}
\end{lstlisting}
